\documentclass[10pt,twocolumn]{article}
\usepackage[T1]{fontenc}
\usepackage[utf8]{inputenc}
\usepackage{lmodern}
\usepackage[margin=1.8cm,top=2.4cm,bottom=2cm,columnsep=0.6cm]{geometry}
\usepackage{fancyhdr}
\usepackage{titlesec}
\usepackage{booktabs}
\usepackage{amsmath,amssymb,amsfonts}
\usepackage{microtype}
\usepackage{xcolor}
\usepackage{mdframed}
\usepackage{parskip}
\usepackage{enumitem}
\usepackage{hyperref}

\definecolor{rulered}{RGB}{180,30,30}
\definecolor{boxgray}{RGB}{245,245,245}
\definecolor{keywordgray}{RGB}{100,100,100}

\pagestyle{fancy}
\fancyhf{}
\fancyhead[L]{\textsc{\small Claude Mag}}
\fancyhead[C]{\small\textit{Machine Learning Research Digest}}
\fancyhead[R]{\small 2026-08-31}
\fancyfoot[C]{\small\thepage}
\renewcommand{\headrulewidth}{0.8pt}

\titleformat{\section}{\normalsize\bfseries\color{rulered}}{}{0pt}{}%
  [\vspace{-4pt}\color{rulered}\rule{\columnwidth}{0.4pt}]
\titlespacing{\section}{0pt}{8pt}{4pt}

\hypersetup{colorlinks=true,urlcolor=rulered,linkcolor=black}

\begin{document}
\setlength{\parindent}{0pt}
\setlength{\parskip}{4pt}

{\small\color{keywordgray}\textbf{Keywords:}
  video generation \textbullet{}
  visual reasoning \textbullet{}
  benchmark \textbullet{}
  diffusion models}\par\vspace{4pt}

{\LARGE\bfseries\raggedright
  VGI-Bench: Probing Visual Intelligence
  in Video Generation Models}\par\vspace{4pt}

{\large\itshape\raggedright
  27 Tasks and 810 Instances Reveal That Even the Strongest
  Video Generator Fails Half the Time on Visually Grounded Reasoning}\par\vspace{6pt}

{\small\textbf{By} Xuan He et al.\ (23 authors)}\par\vspace{2pt}

{\small\color{keywordgray}arXiv:2608.19583 \textbullet{} Preprint, August 2026}
\par\vspace{2pt}\noindent{\color{rulered}\rule{\columnwidth}{0.6pt}}\vspace{6pt}

Video generation models have surpassed human perceptual quality ratings
on many production metrics, but can they \emph{reason visually}
through their generated frames? VGI-Bench probes this question with
27 tasks spanning physics, spatial logic, counting, and rule-following,
each scored by an executable verifier that accepts any feasible
solution. The answer is sobering: Seedance 2.0, the strongest evaluated
model, achieves 51.0\%; denoising trajectory analysis reveals that
models commit to visual hypotheses early and almost never self-correct.

\begin{mdframed}[backgroundcolor=boxgray,linecolor=rulered,linewidth=1pt,
    innerleftmargin=6pt,innerrightmargin=6pt,
    innertopmargin=6pt,innerbottommargin=6pt]
\textbf{\small AT A GLANCE}\par\vspace{4pt}
\begin{description}[leftmargin=0pt,labelsep=4pt,itemsep=2pt,parsep=0pt,
    font=\small\bfseries]
  \item[Problem:] Existing video generation benchmarks reward perceptual
    plausibility rather than visually grounded reasoning; current models
    may look good while failing to satisfy physical or logical constraints.
  \item[Why it matters:] Video generation is increasingly used for
    simulation, planning, and reasoning augmentation; measuring visual
    intelligence is essential to evaluate genuine utility.
  \item[Method:] 27 tasks $\times$ 810 instances with a two-level taxonomy;
    executable verifiers accept any feasible solution (not just ground truth);
    inputs designed for video model visual priors.
  \item[Result:] Seedance 2.0 achieves 51.0\% Verifier Acceptance Rate (VAR);
    most models score below 40\%; denoising analysis shows no self-correction.
  \item[Evidence:] Physical constraint tasks are hardest (38\%);
    failure modes: physical collapse (31\%), rule violation (28\%),
    object inconsistency (24\%).
\end{description}
\end{mdframed}

\section{The Big Picture}

Video generation models are no longer merely artists. They are
increasingly deployed as simulators for planning, world-model probes,
and reasoning-augmented generation — roles that require genuine visual
intelligence, not just surface realism.

Recent work has suggested that video generation models may exhibit
emergent visual reasoning: given a visual prompt depicting an initial
state, can the model generate a video that evolves the scene in a
physically and logically correct manner toward a goal? This would
indicate that the model has learned not just to hallucinate plausible
footage, but to solve visual problems through generation.

VGI-Bench is the first benchmark designed to test this hypothesis
systematically, with tasks that have verifiable solutions and inputs
matched to the visual priors of contemporary video generators.

\section{Why Existing Approaches Fall Short}

\textbf{Plausibility vs.\ correctness:} Existing benchmarks (VBench,
FID/FVD-based metrics) measure perceptual quality — whether generated
videos look realistic. A video can be perfectly realistic while showing
physically incorrect dynamics or failing to satisfy an explicit
constraint. Perceptual metrics cannot measure whether the model solved
the visual problem.

\textbf{Input misalignment:} Benchmarks designed for discriminative
VLMs use image-caption pairs as input, which do not match the
generation-conditioned visual priors of diffusion video models. Inputs
must be designed as visual scenes appropriate for the model's training
distribution.

\textbf{Benchmark saturation:} Benchmarks with simple visual tasks
are rapidly saturated by frontier models, providing no differentiation
above a certain capability threshold.

\section{Core Mechanism}

VGI-Bench specifies tasks via \textbf{hidden structured programs}
that define:
\begin{itemize}[itemsep=2pt,parsep=0pt]
  \item An initial visual scene (rendered as an image prompt)
  \item A goal specification or constraint set
  \item An executable verifier that accepts any feasible solution
\end{itemize}

The verifier is key: it does not compare to a fixed ground-truth
video. Instead it takes the final frame (or frame sequence) and
checks whether the specification is satisfied. This allows diverse
valid solutions and avoids rewarding models that happen to generate
the same pixels as a reference.

The \textbf{two-level taxonomy} crosses:
\begin{description}[leftmargin=0pt,itemsep=2pt,parsep=0pt,font=\small\bfseries]
  \item[Task domain (9):] Physics simulation, spatial arrangement,
    counting, pattern completion, causal reasoning, rule-following,
    logical deduction, navigation, object tracking
  \item[Skill tag (5):] Constraint satisfaction, state tracking,
    multi-step planning, self-correction, multi-object coordination
\end{description}

The 27 tasks are chosen from the $9 \times 5$ crossing to maximise
domain and skill diversity while maintaining calibrated difficulty
(tasks are neither trivial for current models nor completely unsolvable).

\section{Technical Formulation}

Let $\mathcal{P}$ be the space of structured specifications and
$f: \mathcal{P} \rightarrow \mathcal{S}$ a rendering function that
maps a specification to a visual scene. A video model $G$ generates
a video $v = G(f(p))$ conditioned on the rendered scene.

The verifier $\mathcal{V}: v \times p \rightarrow \{0, 1\}$ checks
whether video $v$ satisfies specification $p$. The
\textbf{Verifier Acceptance Rate (VAR)} over $M$ instances is:
\[
  \text{VAR}(G) = \frac{1}{M}\sum_{i=1}^M \mathcal{V}(G(f(p_i)),\, p_i)
\]

The benchmark evaluates each model with a single generation per
instance (no rejection sampling), to measure default generation
quality rather than best-of-$k$ performance.

\section{Learning or Inference Procedure}

No training occurs within the benchmark. Models generate videos
from image prompts using their standard inference pipeline (text
conditioning is minimal — only a brief, non-task-specific caption
is provided to avoid information leakage).

\textbf{Denoising trajectory analysis:} For diffusion-based models,
VGI-Bench provides a methodology for inspecting intermediate denoising
steps using a frozen VLM evaluator at regular intervals (every 10\%
of denoising steps). This reveals whether the model establishes the
correct visual structure early or corrects errors later in denoising.

\section{What the Guarantee Says}

The verifier acceptance criterion is a sufficient condition for
correctness on each task (any feasible solution is accepted), making
VAR a lower bound on true task-solving rate. The paper validates
verifier accuracy against human ratings: verifier agreement with
human binary correctness judgments is 94.2\% across a 200-instance
subset.

\section{Experimental Findings}

\begin{itemize}[itemsep=2pt,parsep=0pt]
  \item \textbf{Seedance 2.0:} 51.0\% VAR (strongest model tested)
  \item \textbf{Wan-2.1, Kling-2.0, others:} Below 40\% VAR
  \item \textbf{Per-domain:} Physics constraints hardest
    (Seedance 2.0: 38\%); counting/arrangement most accessible
    (Seedance 2.0: 67\%)
  \item \textbf{Per-skill:} Multi-step planning tasks show largest
    cross-model spread, confirming they are capability-differentiating
  \item \textbf{Denoising trajectories:} Models set visual structure
    in first $\sim$20\% of denoising steps; later steps add detail
    without correcting structural errors; self-correction rate: $<$4\%
    of initially incorrect states across all models
\end{itemize}

\textbf{Failure mode breakdown:}
\begin{itemize}[itemsep=2pt,parsep=0pt]
  \item Physical collapse (gravity, collision violations): 31\%
  \item Rule violation (incorrect step in evolving system): 28\%
  \item Object/state inconsistency (identity drift): 24\%
  \item Correct terminal state, incorrect trajectory: 17\%
\end{itemize}

\section{Ablations and Interpretation}

\textbf{Text conditioning:} Adding task-specific text prompts (hint
at the constraint) improves VAR by $+7.4$ pp on average, but models
remain far from ceiling, suggesting limited text-to-visual reasoning
integration.

\textbf{Best-of-$k$ generation:} Best-of-8 using the verifier as an
oracle oracle selector improves Seedance 2.0 to 66.2\% VAR, showing
that partial capability exists but is unreliably expressed in a single
generation.

\textbf{Taxonomy analysis:} Skill tag ``constraint satisfaction''
is the weakest across models (mean VAR 32\%); ``counting and arrangement''
is the strongest (mean VAR 51\%). This suggests that spatial counting
and arrangement has been inadvertently trained, while hard constraint
satisfaction has not.

\section*{Reference}

Xuan He et al.
\textbf{VGI-Bench: Probing Visual Intelligence in Video
Generation Models.}
arXiv:2608.19583, August 2026.
\url{https://arxiv.org/abs/2608.19583}

\end{document}
